\ifdefined\CRevisionRound\else\def\CRevisionRound{2}\fi
\documentclass[10pt]{article}
\pdfvariable suppressoptionalinfo 767
\usepackage[a4paper,margin=22mm]{geometry}
\usepackage{amsmath,amssymb,amsthm,booktabs}
\usepackage{fontspec}
\setmainfont{TeX Gyre Pagella}
\setmonofont{Latin Modern Mono}
\usepackage{unicode-math}
\setmathfont{TeX Gyre Pagella Math}
\usepackage[english]{babel}
\babelprovide[import]{chinese-simplified}
\babelfont[chinese-simplified]{rm}[Renderer=Harfbuzz,Language=Default]{Droid Sans Fallback}
\usepackage[hidelinks]{hyperref}
\newtheorem{theorem}{Theorem}
\newtheorem{lemma}{Lemma}
\newtheorem{corollary}{Corollary}
\newcommand{\Gal}{\operatorname{Gal}}
\newcommand{\Aut}{\operatorname{Aut}}
\newcommand{\Fix}{\operatorname{Fix}}
\newcommand{\Per}{\operatorname{Per}}
\title{The Full Generic Binary Tower of $X^2+1$:\\
\ifcase\CRevisionRound New Branch Inertia at Every Height
\or Complete Cycle Indices and Galois-Cover Genus
\else From Ramified Inverse Trees to Vanishing Periodic Density\fi}
\author{HCS-C393 theorem and reproducibility package}
\date{5 September 2026}
\begin{document}
\maketitle
\begin{abstract}
The polynomial $X^2+1$ generates the full automorphism group of its generic
binary preimage tree at every height. A new critical value produces one
sibling transposition; conjugation and restriction then force every
last-level swap, without relying on a finite Galois computation.
\ifnum\CRevisionRound>0
We derive the complete cycle-index recursion, the exact fixed-leaf law,
and the genus of each smooth projective Galois-closure curve. The probability
that a uniform height-$n$ automorphism fixes a leaf is asymptotic to $2/n$,
although its expected fixed-leaf count is exactly one.
\fi
\ifnum\CRevisionRound>1
At each fixed height, explicitly good reductions preserve the full group.
Finite-field Chebotarev then implies that the proportion of periodic points
of the original map modulo primes tends to zero. The order of the two
limits is essential and supplies no uniform-in-height estimate.
\fi
Exact small-tree permutations, critical-value identities and residue-field
computations audit the normalization and its failure boundaries.
These are source-local reconstructions with classical Galois and
periodic-density owners. A transcendental generic basepoint, a Frobenius
cycle and an original-map periodic orbit are not interchangeable, and no
target prime-period or spectral identity is asserted.
\end{abstract}
\noindent\textbf{Keywords:} arboreal Galois group; quadratic dynamics; ramification; wreath product; cycle index; periodic density.
\begin{otherlanguage}{chinese-simplified}
\begin{center}{\large 中文摘要}\end{center}
\begin{sloppypar}
本文对二次映射的泛型逆像树证明每一层 \foreignlanguage{english}{Galois} 群都是完整二叉树自同构群。
每层新增临界值给出一个兄弟交换，再由共轭产生全部末层交换。
\ifnum\CRevisionRound>0
进一步推导所有循环类型的精确递推、固定叶概率及 \foreignlanguage{english}{Galois} 闭包曲线亏格；
固定叶概率渐近为二除以树高，而平均固定叶数始终为一。
\fi
\ifnum\CRevisionRound>1
固定树高之后，有限域切博塔廖夫定理给出像集比例，再按正确的极限顺序
证明模素数周期点比例趋于零。该论证不提供对任意树高一致的误差。
\fi
文中区分泛型基点、数值特化、原映射时间和 \foreignlanguage{english}{Frobenius} 时间，
经典理论明确归属，不把算术源结构写成目标零点对应。
\end{sloppypar}
\noindent 关键词：树状\foreignlanguage{english}{Galois}群；二次动力学；分歧；圈积；循环指标；周期点密度。
\end{otherlanguage}

\section{An intrinsic arithmetic source, not a fitted prime clock}
Let $f(X)=X^2+1$, let $t$ be transcendental over $\mathbb Q$, and let
$K_n$ be the splitting field of $f^n(X)-t$ over $\mathbb Q(t)$.
The rooted tree $T_n$ has root $t$, level-$j$ vertices $f^j(X)=t$,
and edges $x\mapsto f(x)$. Write
\[
 G_n=\Aut(T_n),\qquad d_n=|G_n|=2^{\,2^n-1}.
\]
The generic basepoint is part of the frozen source.
No theorem below says that every specialization $t=a\in\mathbb Q$
has this Galois group.

Odoni \cite{odoni} established the Galois theory of iterated polynomials.
Pink \cite{pink} studies quadratic iterated monodromy with infinite
postcritical orbits; not every quadratic rational map in that broader
class has a full tree group. We give the specific new-branch proof needed
for this polynomial. The wreath-product approach to periodic proportions
belongs to Juul--Kurlberg--Madhu--Tucker \cite{jkmt}.
Our finite audit and explicit source comparison do not certify global
literature novelty.

Within the project, C374 treated the numerical basepoint-two power-map
tower, with an affine index-two image. The present source has full
nonabelian wreath growth, new ramification at every height and an original
finite-field dynamical consequence. It is not another parameter table for
that affine tower. All claims below concern native arithmetic; they do
not insert a prime list or a logarithmic roof into the map.

\section{The complete group from one new branch}
Set $c_0=0$ and $c_i=f^i(0)$, so that
\[
 0,1,2,5,26,\ldots
\]
are distinct: $c_{i+1}=c_i^2+1>c_i$ for $i\ge1$.
The chain rule gives
\begin{equation}\label{der}
 (f^n)'(X)=2^n\prod_{j=0}^{n-1}f^j(X).
\end{equation}
Thus the finite branch values at height $n$ are precisely
$c_1,\ldots,c_n$.
The rational function $f^n$ has degree $2^n$, so
$[\overline{\mathbb Q}(X):\overline{\mathbb Q}(f^n(X))]=2^n$.
Equivalently $f^n(X)-t$ is irreducible over $\overline{\mathbb Q}(t)$,
and geometric monodromy is transitive on its leaves.

\begin{theorem}\label{full}
For every $n\ge1$, both the geometric and arithmetic Galois groups of
$f^n(X)-t$ are $G_n$. The restriction to height $n-1$ is onto with kernel
$(C_2)^{2^{n-1}}$. The infinite tower has group $\Aut(T_\infty)$.
\end{theorem}
\begin{proof}
At $t=c_n$ the only critical point in the fibre is $X=0$.
Indeed $f^j(X)=0$ for $j\ge1$ would imply
$f^n(X)=c_{n-j}\ne c_n$. Near zero,
\[
 f^n(X)=c_n+A_nX^2+O(X^4),\qquad
 A_n=2^{n-1}\prod_{i=1}^{n-1}c_i\ne0.
\]
Local inertia is therefore a single transposition of two leaves.
They are siblings because their common image under $f$ is $1$.
The previous-height cover is unramified at this new branch value, so
the transposition belongs to the last-level restriction kernel.

Inductively the group at height $n-1$ is $G_{n-1}$; restriction between
splitting fields is surjective. The new kernel is contained in the group
of independent swaps of the $2^{n-1}$ sibling pairs.
Conjugate the one new swap by elements transitive on their parents.
Every sibling swap is obtained, and these generate the entire kernel.
The degree-two case starts the induction.

The arithmetic group contains the full geometric group and embeds in
$G_n$, so it too is full. Thus the constant field is $\mathbb Q$.
Compatible surjective restrictions give the full inverse limit:
an automorphism of the union is exactly a family of compatible
finite-level automorphisms.
\end{proof}
This proof derives the kernel rather than inferring it from a group-order
fit. It also records why a postcritical collision would break the argument.

\ifnum\CRevisionRound>0
\section{All branch inertia and the Galois-closure genus}
At $t=c_i$, the ramified points in the degree-$2^n$ covering are the
$2^{n-i}$ roots of $f^{n-i}(X)=0$. They are simple because no earlier
critical value is zero. Each has ramification index two.
The finite inertia permutation therefore has
\[
 2^{n-i}\text{ two-cycles},\qquad
 2^n-2^{n-i+1}\text{ fixed leaves}.
\]
At infinity the polynomial has a single point of index $2^n$, with
inertia a $2^n$-cycle. All inertia in the Galois closure is tame and
there are no further branch values.

\begin{theorem}\label{genus}
The smooth projective Galois-closure curve of $K_n/\mathbb Q(t)$ has genus
\begin{equation}\label{gen}
 g_n=1+\frac{d_n}{2}\left(\frac n2-1-2^{-n}\right).
\end{equation}
\end{theorem}
\begin{proof}
For a tame Galois cover of degree $d_n$, each base branch point of inertia
order $e$ contributes $d_n(1-1/e)$ to Riemann--Hurwitz.
There are $n$ finite branch values with $e=2$ and one at infinity with
$e=2^n$. Consequently
\[
 2g_n-2=d_n[-2+n/2+(1-2^{-n})],
\]
which is \eqref{gen}.
\end{proof}
For example $g_1=g_2=0$ and $g_3=25$.
This is not the genus of the original polynomial-cover source,
which is the projective line.

\section{The full cycle index and fixed-leaf law}
Let $Z_n(x_1,x_2,\ldots)$ be the mean cycle monomial of a uniform element
of $G_n$, with $Z_0=x_1$.
The wreath decomposition gives the complete recursion
\begin{equation}\label{index}
 Z_n(x)=\frac12\left[Z_{n-1}(x)^2+
                 Z_{n-1}(x_2,x_4,\ldots)\right].
\end{equation}
Without the top swap, the two independent cycle partitions add.
With the swap, two steps restrict to the product of the two independent
subtree permutations, which is uniform; every resulting cycle length doubles.
This proves every coefficient, not only a fixed-root marginal.

For an element with $m_l$ cycles of length $l$,
\[
 \#\Fix(g^r)=\sum_{l\mid r}lm_l,\qquad
 \det(I-uU_g)=\prod_l(1-u^l)^{m_l}.
\]
Here $U_g$ is the native finite leaf permutation. Frobenius repetition
changes that permutation; it is neither tree height nor the original
map's iterate period. The finite determinant is not a target Fredholm operator.

\begin{theorem}\label{fpp}
If $\delta_n$ is the probability of at least one fixed leaf, then
\[
 \delta_0=1,\qquad
 \delta_n=\delta_{n-1}-\tfrac12\delta_{n-1}^2,\qquad
 n\delta_n\longrightarrow2.
\]
The mean fixed-leaf count is one at every height.
\end{theorem}
\begin{proof}
A top swap fixes no leaf. Without it, a fixed leaf exists unless both
independent subtrees have none. This proves the recursion.
The positive sequence decreases to zero, since its limit satisfies
$L=L-L^2/2$. Also
\[
 \frac1{\delta_n}-\frac1{\delta_{n-1}}
   =\frac1{2(1-\delta_{n-1}/2)}\longrightarrow\frac12 .
\]
Cesàro summation proves $n\delta_n\to2$. The expected count equals one
by transitivity and the orbit-counting lemma, or directly by
\eqref{index}.
\end{proof}
Thus a vanishing chance of any fixed leaf coexists with constant mean count.
The distinction is intrinsic to the full group and cannot be read from one
sample Frobenius element.
\fi

\ifnum\CRevisionRound>1
\section{Good reductions and the ordered two-limit argument}
For fixed $n$ define the explicit nonzero integer
\begin{equation}\label{bad}
 B_n=2\prod_{0\le i<j\le n}(c_j-c_i).
\end{equation}
If $p\nmid B_n$, the critical values through height $n$, including zero,
remain distinct modulo $p$. The same local proof of Theorem \ref{full}
works over $\overline{\mathbb F}_p(t)$. All inertia orders are powers of
two, so the reduction is tame. Its geometric group is full; its
arithmetic group contains that group and embeds in it, so the cover
is regular over $\mathbb F_p$.

We explicitly import the classical finite-field Chebotarev theorem for
regular covers, obtained from the Weil bounds on twists.
For a fixed height and a union of conjugacy classes, the proportion of
unramified rational basepoints has the group-theoretic proportion plus
$O_n(p^{-1/2})$. This is the same input used in \cite{jkmt},
Section 5 and Proposition 5.3. The implicit constant may depend on the
large degree and genus in \eqref{gen}; it is not uniform in height.

At an unramified $a\in\mathbb F_p$, the factor degrees of
$f^n(X)-a$ are exactly the cycle lengths of its Frobenius leaf permutation.
Thus every cycle-index coefficient in \eqref{index} has that fixed-height
distributional interpretation. In particular, a root in $\mathbb F_p$
exists precisely when Frobenius fixes a leaf.
At most $n$ finite branch values are excluded, so
\begin{equation}\label{image}
 \frac{|f^n(\mathbb F_p)|}{p}=\delta_n+O_n(p^{-1/2}).
\end{equation}

\begin{corollary}\label{density}
For the original map $f(x)=x^2+1$ on $\mathbb F_p$,
\[
 \lim_{\substack{p\to\infty\\p\ {\rm prime}}}
        \frac{|\Per(f,\mathbb F_p)|}{p}=0.
\]
\end{corollary}
\begin{proof}
A periodic point has $n$ backward steps along its cycle for every $n$.
Hence $\Per(f,\mathbb F_p)\subset f^n(\mathbb F_p)$.
For each fixed $n$, only finitely many primes divide \eqref{bad};
taking the prime limsup in \eqref{image} gives an upper bound $\delta_n$.
Only now send $n$ to infinity and apply Theorem \ref{fpp}.
Nonnegativity gives the zero limit.
\end{proof}
This is a source specialization of the classical wreath-product density
mechanism \cite{jkmt}, with the generic-group and limit-order hypotheses
made explicit. It supplies no rate in $p$ and makes no assertion with
a prime-dependent cutoff $n(p)$.

\section{Countercontrols, independent evidence and the route verdict}
The simpler source $X^2$ already has a critical collision at height one
and cyclic geometric monodromy. The neighboring source $X^2-1$ has
critical orbit $0,-1,0$, so the new-branch induction fails.
Reduction at bad primes supplies an intrinsic third control.
Indeed no finite residue field can contain infinitely many distinct
critical values. There is no prime that satisfies \eqref{bad} at every
height. Composite moduli are not fields and are not called single Frobenius owners.

The exact ledger contains eight tower rows, cycle indices through height
five, thirteen fixed-leaf probabilities, and 150 map-level rows over the
25 primes below 100. The independent checker directly enumerates 32,907
tree permutations through height four, then forms the next class convolution.
It computes original-map periodic points by functional-graph path detection,
independently of the producer's stabilized image sets.
SymPy checks 14 identities and independently factors 632 residue polynomials.
An initial installed-backend sorting error aborted that lane; selecting
Python ground types before importing SymPy repaired the executable.
It did not supply a mathematical counterexample or change the evidence.

These finite receipts check signs, multiplicities, bad-prime exclusions and
schema handling. They do not prove the infinite group or the density limit.
Tree height, original-map iteration and Frobenius iteration remain distinct.
The source supplies genuine arithmetic ramification, but no rational prime
has been identified with an original orbit of period $\log p$.

The conservative arithmetic grade is
\[
\mathrm{A0\_STRUCTURAL\_ARITHMETIC\_RELATION}.
\]
The remaining grades are
$(\mathrm{A1\_WEAK},\mathrm{A2\_FAIL},\mathrm{A3\_FAIL},
\mathrm{A4\_FORMAL\_HINT})$, overall exploratory.
Finite leaf unitaries are not a global Hilbert--Pólya construction.
\par\noindent The firewall \texttt{NO\_BAD\_EULER\_OR\_ROOT\_NUMBER} remains in force,
with every target/Route-B flag false.
\fi

\section{Conclusion}
New critical inertia closes the entire generic binary Galois tower.
\ifnum\CRevisionRound>0
Its cycle indices and Galois-cover genus follow at every height, not only
on the finite verification grid.
\fi
\ifnum\CRevisionRound>1
Fixed-level arithmetic reduction then reaches an original finite-field
periodic-density conclusion through an ordered, nonuniform two-limit argument.
\fi
This is one paper, not three publications; its present record is
Round \ifcase\CRevisionRound zero\or one\else two\fi.
The accompanying proof ledger, independent checks and source audit delimit
the complete source theorem from the unachieved target bridge.

\begin{thebibliography}{9}
\bibitem{odoni}
R. W. K. Odoni.
The Galois Theory of Iterates and Composites of Polynomials.
\emph{Proceedings of the London Mathematical Society}
s3-51(3), 385--414 (1985).
\href{https://doi.org/10.1112/plms/s3-51.3.385}{doi:10.1112/plms/s3-51.3.385}.
\bibitem{pink}
R. Pink.
Profinite iterated monodromy groups arising from quadratic morphisms
with infinite postcritical orbits.
Preprint (2013), \href{https://arxiv.org/abs/1309.5804}{arXiv:1309.5804}.
\bibitem{jkmt}
J. Juul, P. Kurlberg, K. Madhu and T. J. Tucker.
Wreath Products and Proportions of Periodic Points.
\emph{International Mathematics Research Notices}
2016(13), 3944--3969 (2016; online 2015).
\href{https://doi.org/10.1093/imrn/rnv273}{doi:10.1093/imrn/rnv273}.
Accessed proof version: \href{https://arxiv.org/abs/1410.3378}{arXiv:1410.3378}.
\end{thebibliography}
\end{document}
